% =====================================================================
% Übungskapitel zu Kapitel 3 (Stochastik) – steht direkt hinter Stochastik.tex
% Nummerierung: Ü 3.n ; Lösungen folgen in einem späteren Schritt.
% =====================================================================
\Uebungskapitel{Übungen zur Stochastik}

\Uebungsabschnitt{Wahrscheinlichkeitsbegriffe und bedingte Wahrscheinlichkeit}

\begin{uebung}*{Ereignisse und Axiome}
Ein Laplace-Würfel wird einmal geworfen. Betrachtet werden die Ereignisse $A$: \glqq Die Augenzahl ist gerade.\grqq{} und $B$: \glqq Die Augenzahl ist größer als 4.\grqq
\begin{enumerate}[label=\alph*)]
  \item Geben Sie den Ergebnisraum $\Omega$ sowie die Ereignisse $A$, $B$, $\bar{A}$, $A\cap B$ und $A\cup B$ in Mengenschreibweise an.
  \item Berechnen Sie $P(A\cap B)$ und $P(A\cup B)$.
  \item Entscheiden Sie, ob die Ereignisse $A$ und $B$ unvereinbar sind, und begründen Sie Ihre Entscheidung.
  \item Begründen Sie mithilfe der Grundaxiome der Stochastik, dass für jedes Ereignis $E$ gilt: $P(\bar{E}) = 1-P(E)$.
\end{enumerate}
\end{uebung}

\begin{uebung}{Bedingte Wahrscheinlichkeit}
Eine Krankheit tritt bei $2\,\%$ der Bevölkerung auf. Ein Schnelltest zeigt bei einer erkrankten Person mit einer Wahrscheinlichkeit von $95\,\%$ ein positives Ergebnis. Bei einer nicht erkrankten Person zeigt er mit einer Wahrscheinlichkeit von $10\,\%$ fälschlicherweise ebenfalls ein positives Ergebnis. Eine zufällig ausgewählte Person wird getestet.\\
Betrachtet werden die Ereignisse $K$: \glqq Die Person ist erkrankt.\grqq{} und $T$: \glqq Der Test ist positiv.\grqq
\begin{enumerate}[label=\alph*)]
  \item Stellen Sie den Sachverhalt in einem vollständig beschrifteten Baumdiagramm dar.
  \item Berechnen Sie die Wahrscheinlichkeit dafür, dass der Test positiv ausfällt.
  \item Berechnen Sie $P_T(K)$ und interpretieren Sie das Ergebnis im Sachzusammenhang.
  \item Beurteilen Sie anhand Ihres Ergebnisses aus c) die Aussagekraft eines positiven Testergebnisses.
\end{enumerate}
\end{uebung}

\Uebungsabschnitt{Zufallsgrößen, Erwartungswert und Varianz}

\begin{uebung}*{Glücksrad}
Ein Glücksrad besteht aus drei Sektoren. Der Zeiger bleibt mit der Wahrscheinlichkeit $\frac{1}{2}$ im Sektor \glqq 0\,€\grqq, mit der Wahrscheinlichkeit $\frac{1}{3}$ im Sektor \glqq 1\,€\grqq{} und mit der Wahrscheinlichkeit $\frac{1}{6}$ im Sektor \glqq 4\,€\grqq{} stehen. Ausgezahlt wird der Betrag im jeweiligen Sektor. Der Einsatz für ein Spiel beträgt $1$\,€. Die Zufallsgröße $X$ beschreibt den Gewinn des Spielers (Auszahlung minus Einsatz) in Euro.
\begin{enumerate}[label=\alph*)]
  \item Geben Sie die Wahrscheinlichkeitsverteilung von $X$ in einer Tabelle an.
  \item Zeigen Sie, dass das Spiel fair ist.
  \item Berechnen Sie die Varianz und die Standardabweichung von $X$.
  \item Geben Sie den Wert $F(0) = P(X\leq 0)$ der kumulativen Verteilungsfunktion von $X$ an.
  \item Der Betreiber möchte im Mittel $0{,}25$\,€ pro Spiel einnehmen. Dazu soll nur der Betrag im Sektor \glqq 4\,€\grqq{} geändert werden. Bestimmen Sie den neuen Betrag.
\end{enumerate}
\end{uebung}

\Uebungsabschnitt{Kombinatorik und hypergeometrische Verteilung}

\begin{uebung}*{Kombinatorik}
Geben Sie jeweils einen Term für die Anzahl der Möglichkeiten an und berechnen Sie diese, falls möglich, ohne Taschenrechner.
\begin{enumerate}[label=\alph*)]
  \item Ein Zahlenschloss hat vier Ringe mit den Ziffern $0$ bis $9$.
  \item Ein vierstelliger Code wird aus den Ziffern $0$ bis $9$ gebildet, wobei keine Ziffer mehrfach vorkommen darf.
  \item Fünf verschiedene Bücher werden nebeneinander in ein Regal gestellt.
  \item Aus fünf Personen wird ein Zweierteam ausgewählt.
  \item Beim Lotto \glqq 6 aus 49\grqq{} werden sechs von 49 Kugeln ohne Zurücklegen und ohne Beachtung der Reihenfolge gezogen. (Nur Term angeben.)
\end{enumerate}
\end{uebung}

\begin{uebung}{Hypergeometrische Verteilung}
In einer Lieferung von $20$ Glühbirnen sind $4$ defekt. Es werden $5$ Glühbirnen auf einmal (ohne Zurücklegen) entnommen. Die Zufallsgröße $X$ gibt die Anzahl der defekten Glühbirnen in der Stichprobe an.
\begin{enumerate}[label=\alph*)]
  \item Begründen Sie, dass $X$ hypergeometrisch verteilt ist.
  \item Berechnen Sie die Wahrscheinlichkeit dafür, dass genau eine defekte Glühbirne entnommen wird.
  \item Berechnen Sie die Wahrscheinlichkeit dafür, dass mindestens eine defekte Glühbirne entnommen wird.
\end{enumerate}
\end{uebung}

\Uebungsabschnitt{Bernoulli-Kette und Binomialverteilung}

\begin{uebung}{Multiple-Choice-Test}
Ein Test besteht aus $10$ Fragen mit jeweils vier Antwortmöglichkeiten, von denen genau eine richtig ist. Eine Person kreuzt bei jeder Frage zufällig genau eine Antwort an. Die Zufallsgröße $X$ gibt die Anzahl der richtig beantworteten Fragen an.
\begin{enumerate}[label=\alph*)]
  \item Begründen Sie, dass $X$ als binomialverteilt angesehen werden kann, und geben Sie die Parameter an.
  \item Berechnen Sie die Wahrscheinlichkeiten folgender Ereignisse:\\
  $A$: \glqq Genau drei Antworten sind richtig.\grqq\\
  $B$: \glqq Höchstens zwei Antworten sind richtig.\grqq\\
  $C$: \glqq Der Test ist bestanden, d.\,h. mindestens fünf Antworten sind richtig.\grqq\\
  $D$: \glqq Mindestens zwei, aber höchstens vier Antworten sind richtig.\grqq
  \item Berechnen Sie den Erwartungswert und die Standardabweichung von $X$ und geben Sie die Anzahl richtiger Antworten mit der größten Wahrscheinlichkeit an.
\end{enumerate}
\end{uebung}

\begin{uebung}{3-m-Aufgaben}
\begin{enumerate}[label=\alph*)]
  \item Ein Bauteil ist mit einer Wahrscheinlichkeit von $5\,\%$ defekt. Bestimmen Sie, wie viele Bauteile man mindestens prüfen muss, damit sich darunter mit einer Wahrscheinlichkeit von mindestens $95\,\%$ mindestens ein defektes Bauteil befindet.
  \item Bestimmen Sie, wie groß die Wahrscheinlichkeit $p$ für einen Treffer mindestens sein muss, damit bei $20$ Versuchen mit einer Wahrscheinlichkeit von mindestens $90\,\%$ mindestens ein Treffer erzielt wird.
\end{enumerate}
\end{uebung}

\begin{uebung}{\texorpdfstring{\(\sigma\)}{Sigma}-Umgebung}
Die Zufallsgröße $X$ ist nach $B(100;0{,}3)$ verteilt.
\begin{enumerate}[label=\alph*)]
  \item Berechnen Sie den Erwartungswert $\mu$ und die Standardabweichung $\sigma$ von $X$.
  \item Bestimmen Sie die $2\sigma$-Umgebung von $\mu$ und geben Sie alle ganzzahligen Werte von $X$ an, die in ihr liegen.
  \item Berechnen Sie die Wahrscheinlichkeit dafür, dass $X$ einen Wert aus dieser Umgebung annimmt.
  \item Bei einer Durchführung des Zufallsexperiments werden $45$ Treffer erzielt. Beurteilen Sie dieses Ergebnis.
\end{enumerate}
\end{uebung}

\Uebungsabschnitt{Hypothesentest}

\begin{uebung}{Fehler beim Testen}
Ein Hersteller behauptet, dass höchstens $10\,\%$ seiner Bauteile defekt sind. Ein Abnehmer möchte diese Behauptung überprüfen und testet die Nullhypothese $H_0$: \glqq Der Anteil defekter Bauteile beträgt höchstens $10\,\%$.\grqq{} gegen die Gegenhypothese $H_1$: \glqq Der Anteil defekter Bauteile beträgt mehr als $10\,\%$.\grqq{} Dazu werden $50$ zufällig ausgewählte Bauteile geprüft. Die Testgröße $X$ ist die Anzahl der defekten Bauteile in der Stichprobe. Der Abnehmer legt zunächst folgende Entscheidungsregel fest: Sind mindestens $9$ Bauteile defekt, wird $H_0$ abgelehnt.
\begin{enumerate}[label=\alph*)]
  \item Geben Sie den kritischen Bereich $K$ und den Annahmebereich $\bar{K}$ an.
  \item Beschreiben Sie den Fehler erster Art und den Fehler zweiter Art im Sachzusammenhang.
  \item Berechnen Sie die Wahrscheinlichkeit für den Fehler erster Art, wenn der Anteil defekter Bauteile genau $10\,\%$ beträgt.
  \item Tatsächlich sind $20\,\%$ der Bauteile defekt. Berechnen Sie für diesen Fall die Wahrscheinlichkeit für den Fehler zweiter Art.
  \item Der Abnehmer möchte, dass die Wahrscheinlichkeit für den Fehler erster Art (bei einem Anteil von genau $10\,\%$) höchstens $5\,\%$ beträgt. Bestimmen Sie den kleinsten Wert von $k$, für den der kritische Bereich $K = \{k; k+1; \ldots; 50\}$ diese Bedingung erfüllt.
  \item Beschreiben Sie, wie sich die Änderung der Entscheidungsregel aus e) auf die Wahrscheinlichkeit für den Fehler zweiter Art auswirkt, und nennen Sie eine Möglichkeit, beide Fehlerwahrscheinlichkeiten zu verringern.
\end{enumerate}
\end{uebung}

\Uebungsabschnitt{Stetige Zufallsgrößen und Normalverteilung}

\begin{uebung}*{Dichtefunktion}
Gegeben ist die Funktion $f$ mit
$$f(x)= \begin{cases}
              k\cdot x\cdot(4-x) & \text{für}\hspace{0.5cm} 0\leq x\leq 4\\
              0 & \text{sonst}
         \end{cases}\qquad k\in\mathds{R}^+.$$
\begin{enumerate}[label=\alph*)]
  \item Bestimmen Sie $k$ so, dass $f$ eine Dichtefunktion ist.
  \item Die stetige Zufallsgröße $X$ besitzt die Dichtefunktion $f$. Berechnen Sie $P(1\leq X\leq 3)$.
  \item Begründen Sie ohne Rechnung mithilfe einer Symmetrieüberlegung, dass $E(X)=2$ gilt.
  \item Berechnen Sie die Varianz von $X$.
  \item Geben Sie $P(X=2)$ an und begründen Sie Ihre Angabe.
\end{enumerate}
\end{uebung}

\begin{uebung}{Gauß'sche Glockenkurve}
\begin{enumerate}[label=\alph*)]
  \item Gegeben ist die Funktion $\varphi_{0;1}$ mit \[\varphi_{0;1}(x) = \dfrac{1}{\sqrt{2\pi}}\cdot e^{-\frac{x^2}{2}}.\] Zeigen Sie ohne Taschenrechner, dass ihr Graph an den Stellen $x=-1$ und $x=1$ Wendepunkte besitzt.
  \item Der Hochpunkt des Graphen einer Gauß'schen Glockenfunktion $\varphi_{\mu;\sigma}$ ist $H(10|0{,}2)$. Bestimmen Sie $\mu$ und $\sigma$ (auf eine Nachkommastelle gerundet) und geben Sie die Wendestellen des Graphen an.
  \item Beschreiben Sie, wie sich der Graph von $\varphi_{\mu;\sigma}$ verändert, wenn $\sigma$ bei gleichem $\mu$ verdoppelt wird.
\end{enumerate}
\end{uebung}

\begin{uebung}{Normalverteilung und \texorpdfstring{\(\sigma\)}{Sigma}-Regeln}
Die Füllmenge von Saftflaschen (in ml) wird durch eine $N_{500;4}$-verteilte Zufallsgröße $X$ beschrieben.
\begin{enumerate}[label=\alph*)]
  \item Bestimmen Sie mithilfe der $\sigma$-Regeln näherungsweise die Wahrscheinlichkeit dafür, dass eine zufällig ausgewählte Flasche weniger als $492$\,ml enthält.
  \item Berechnen Sie die Wahrscheinlichkeit dafür, dass die Füllmenge zwischen $495$\,ml und $505$\,ml liegt.
  \item Bestimmen Sie ein zu $\mu$ symmetrisches Intervall, in dem die Füllmenge mit einer Wahrscheinlichkeit von etwa $95\,\%$ liegt.
  \item Es werden $20$ Flaschen zufällig ausgewählt. Berechnen Sie die Wahrscheinlichkeit dafür, dass mindestens eine davon weniger als $492$\,ml enthält. Verwenden Sie dazu den exakten Wert aus der Normalverteilung.
  \item Beurteilen Sie, ob die Modellierung der Füllmenge durch eine Normalverteilung sinnvoll ist.
\end{enumerate}
\end{uebung}

\UebungskapitelEnde
